\chapter{Przestrzenie odwzorowań i topologia zwarto-otwarta}
\label{ch:topog-przestrzenie-odwzorowan}

\noindent\textbf{Cel rozdziału.}
Wcześniej rozpatrywaliśmy ciągłość pojedynczego
odwzorowania $X\to Y$.
Teraz sam zbiór wszystkich takich odwzorowań
uczynimy przestrzenią topologiczną.
Pozwoli to badać rodziny odwzorowań, ich granice,
odwzorowanie ewaluacji i zależność funkcji
od parametrów. Szczególną rolę będzie odgrywać
\emph{zbieżność jednostajna na zbiorach zwartych}.
Zwartość lokalna z rozdziału~14 będzie potrzebna
dopiero przy ewaluacji i prawie wykładniczym;
sama zgodność dwóch podstawowych topologii
na zbiorze funkcji ciągłych nie wymaga
tego założenia.

\section{Topologia zwarto-otwarta}
\label{sec:topog-odwzorowania-zwarto-otwarta}

\begin{definicja}[Przestrzeń odwzorowań ciągłych]
\label{def:topog-odwzorowania-Ck}
Niech $X,Y$ będą przestrzeniami topologicznymi.
Oznaczmy przez $C(X,Y)$ zbiór wszystkich
odwzorowań ciągłych $f\colon X\to Y$.
Dla każdego zwartego $K\subseteq X$
i każdego otwartego $U\subseteq Y$
definiujemy
\[
[K,U]=\{f\in C(X,Y):f(K)\subseteq U\}.
\]
\emph{Topologią zwarto-otwartą}
nazywamy topologię na $C(X,Y)$
generowaną przez rodzinę wszystkich
zbiorów $[K,U]$ jako podbazę.
Przestrzeń $C(X,Y)$ z tą topologią
oznaczamy przez $C_k(X,Y)$.
\end{definicja}

\begin{uwaga}[Dlaczego podbaza?]
Zbiory $[K,U]$ nie muszą same
stanowić bazy: warunek jednoczesnego
spełnienia $f(K_1)\subseteq U_1$
i $f(K_2)\subseteq U_2$ wymaga
przecięcia dwóch zbiorów podbazowych.
Bazę tworzą zatem ich przecięcia
skończone, a topologię --- dowolne sumy
takich przecięć. Dla $K=\varnothing$
mamy $[K,U]=C(X,Y)$.
\end{uwaga}

\begin{stwierdzenie}[Oddzielanie punktów i zbieżność punktowa]
\label{prop:topog-odwzorowania-oddzielanie}
Niech $X,Y$ będą przestrzeniami topologicznymi.
Topologia zwarto-otwarta na $C(X,Y)$
jest co najmniej tak silna jak topologia
zbieżności punktowej, czyli topologia
podprzestrzeni odziedziczona z iloczynu $Y^X$.
Jeżeli $Y$ jest Hausdorffa,
to $C_k(X,Y)$ jest Hausdorffa.
\end{stwierdzenie}
\begin{proof}
Dla $x\in X$ singleton $\{x\}$ jest
zwarty bez dodatkowych założeń na $X$.
Przeciwobraz zbioru otwartego
$U\subseteq Y$ przez ewaluację
$f\mapsto f(x)$ ma postać
$[\{x\},U]$, a więc jest otwarty
w topologii zwarto-otwartej.
Ponieważ topologia punktowa
jest początkowa względem wszystkich
tych ewaluacji, jest nie silniejsza.

Niech teraz $f,g\in C(X,Y)$ będą różne.
Istnieje $x\in X$ z $f(x)\ne g(x)$.
Wybierzmy rozłączne otwarte
$U,V\subseteq Y$ zawierające
odpowiednio $f(x)$ i $g(x)$.
Wtedy $[\{x\},U]$ oraz $[\{x\},V]$
są rozłącznymi otoczeniami $f$ i $g$.
\end{proof}

\begin{stwierdzenie}[Składanie z ustalonym odwzorowaniem]
\label{prop:topog-odwzorowania-skladanie}
Niech $X,X',Y,Z$ będą przestrzeniami
topologicznymi.
Jeśli $u\colon X'\to X$ i
$v\colon Y\to Z$ są ciągłe, to
odwzorowania
\[
u^\ast\colon C_k(X,Y)\to C_k(X',Y),
\quad u^\ast(f)=f\circ u,
\]
oraz
\[
v_\ast\colon C_k(X,Y)\to C_k(X,Z),
\quad v_\ast(f)=v\circ f
\]
są ciągłe.
\end{stwierdzenie}
\begin{proof}
Dla zwartego $K'\subseteq X'$
i otwartego $U\subseteq Y$ mamy
\[
(u^\ast)^{-1}([K',U])=[u(K'),U].
\]
Zbiór $u(K')$ jest zwarty
jako obraz ciągły zwartego $K'$,
więc prawa strona jest otwarta.
Dla zwartego $K\subseteq X$
i otwartego $V\subseteq Z$ mamy
\[
(v_\ast)^{-1}([K,V])
=[K,v^{-1}(V)].
\]
Drugi zbiór jest otwarty,
ponieważ $v^{-1}(V)$ jest otwarty.
Ciągłość obu odwzorowań wynika
z kryterium podbazowego.
\end{proof}

\section{Zbieżność jednostajna na zbiorach zwartych}
\label{sec:topog-odwzorowania-jednostajna-zwarte}

\begin{definicja}[Otoczenia jednostajne na zbiorze zwartym]
\label{def:topog-odwzorowania-jednostajna-zwarte}
Niech $X$ będzie przestrzenią topologiczną,
a $(Y,d)$ przestrzenią metryczną.
Dla $f\in C(X,Y)$, zwartego
$K\subseteq X$ i $\varepsilon>0$ określamy
\[
B_K(f,\varepsilon)=
\{g\in C(X,Y):
 d(f(x),g(x))<\varepsilon
 \text{ dla każdego }x\in K\}.
\]
Rodzina tych zbiorów, przy ustalonym $f$
i wszystkich zwartych $K$ oraz
$\varepsilon>0$, tworzy bazę
otoczeń pewnej topologii na $C(X,Y)$.
Nazywamy ją \emph{topologią zbieżności
jednostajnej na zbiorach zwartych}.
Dla $K=\varnothing$ przyjmujemy
$B_K(f,\varepsilon)=C(X,Y)$.
\end{definicja}
\begin{proof}
Każdy taki zbiór zawiera $f$.
Przecięcie dwóch otoczeń
$B_{K_1}(f,\varepsilon_1)$ i
$B_{K_2}(f,\varepsilon_2)$ zawiera
$B_{K_1\cup K_2}(f,\min\{\varepsilon_1,
\varepsilon_2\})$, gdyż skończona suma
zbiorów zwartych jest zwarta.
Jeśli $g\in B_K(f,\varepsilon)$
i $K\ne\varnothing$, to ciągła
funkcja $x\mapsto d(f(x),g(x))$
osiąga na zwartym $K$ maksimum
$m<\varepsilon$. Dla
$\delta=\varepsilon-m>0$
mamy
$B_K(g,\delta)\subseteq
B_K(f,\varepsilon)$
na mocy nierówności trójkąta.
Dla $K=\varnothing$ własność
jest oczywista. Spełnione są
więc aksjomaty bazy otoczeń.
\end{proof}

\begin{twierdzenie}[Równoważność dwóch topologii]
\label{thm:topog-odwzorowania-zwarto-otwarta-jednostajna}
Niech $X$ będzie dowolną przestrzenią
topologiczną, a $(Y,d)$ przestrzenią
metryczną. Na $C(X,Y)$ topologia
zwarto-otwarta jest równa topologii
zbieżności jednostajnej na zbiorach
zwartych. Nie jest wymagane, aby $X$
była Hausdorffa ani lokalnie zwarta.
\end{twierdzenie}
\begin{proof}
Najpierw niech $f\in[K,U]$,
gdzie $K\subseteq X$ jest zwarty,
a $U\subseteq Y$ jest otwarty.
Jeśli $K=\varnothing$ lub $U=Y$,
to $[K,U]=C(X,Y)$.
W pozostałym przypadku
$D=Y\setminus U$ jest niepusty
i domknięty, a $f(K)$ jest
zawartym w $U$ zbiorem zwartym.
Funkcja
$y\mapsto d(y,D)=\inf_{z\in D}d(y,z)$
jest ciągła, a dla każdego
$y\in f(K)$ ma wartość dodatnią.
Osiąga więc na $f(K)$ dodatnie
minimum $\eta>0$.
Jeśli $g\in B_K(f,\eta/2)$,
to $g(x)\notin D$ dla każdego
$x\in K$, a więc $g(K)\subseteq U$.
Zatem
$B_K(f,\eta/2)\subseteq[K,U]$.
Każdy zbiór podbazowy topologii
zwarto-otwartej jest otwarty
w topologii zbieżności jednostajnej
na zbiorach zwartych.

Odwrotnie, niech
$g\in B_K(f,\varepsilon)$.
Jeśli $K=\varnothing$, nie ma
czego dowodzić. Zwartość $K$
i ciągłość $x\mapsto d(f(x),g(x))$
dają liczbę
$m=\max_{x\in K}d(f(x),g(x))
<\varepsilon$.
Połóżmy $\delta=\varepsilon-m>0$.
Zwartość $g(K)$ zapewnia
istnienie punktów
$y_1,\ldots,y_s\in g(K)$
takich, że kule
$B_d(y_j,\delta/4)$ pokrywają
$g(K)$. Definiujemy zwarte zbiory
\[
K_j=K\cap
g^{-1}(\overline B_d(y_j,\delta/4)),
\qquad j=1,\ldots,s.
\]
Są zwarte, ponieważ są domkniętymi
podzbiorami zwartego $K$,
a ponadto $\bigcup_{j=1}^sK_j=K$.
Rozpatrzmy otwarte w topologii
zwarto-otwartej otoczenie
\[
\mathcal W=
\bigcap_{j=1}^s
[K_j,B_d(y_j,\delta/2)].
\]
Jest to rzeczywiście otoczenie
funkcji $g$: jeśli $x\in K_j$,
to $d(g(x),y_j)\le\delta/4<\delta/2$.
Dla dowolnego $h\in\mathcal W$
i $x\in K_j$ otrzymujemy
\[
\begin{split}
d(f(x),h(x))
&\le d(f(x),g(x))
+d(g(x),y_j)+d(y_j,h(x))\\
&<m+\delta/4+\delta/2
<\varepsilon.
\end{split}
\]
Każdy $x\in K$ należy do
pewnego $K_j$, więc
$\mathcal W\subseteq B_K(f,\varepsilon)$.
Dla każdego $g$ z tego otoczenia
znaleźliśmy mniejsze otoczenie
zwarto-otwarte; stąd
$B_K(f,\varepsilon)$ jest otwarty
w topologii zwarto-otwartej.
\end{proof}

\begin{wniosek}[Charakterystyka zbieżności sieci]
\label{cor:topog-odwzorowania-sieci}
Niech $X$ będzie przestrzenią topologiczną,
$(Y,d)$ przestrzenią metryczną,
a $(f_\lambda)_{\lambda\in\Lambda}$
siecią w $C(X,Y)$.
Sieć ta jest zbieżna do $f\in C(X,Y)$
w topologii zwarto-otwartej
wtedy i tylko wtedy, gdy dla każdego
zwartego $K\subseteq X$ i każdej
$\varepsilon>0$ istnieje
$\lambda_0\in\Lambda$ takie, że
dla $\lambda\ge\lambda_0$
oraz wszystkich $x\in K$ mamy
$d(f_\lambda(x),f(x))<\varepsilon$.
\end{wniosek}
\begin{proof}
Na mocy poprzedniego twierdzenia
zbiory $B_K(f,\varepsilon)$ stanowią
bazę otoczeń punktu $f$.
Definicja zbieżności sieci daje
dokładnie sformułowany warunek.
\end{proof}

\begin{przyklad}[Zbieżność na zwartych zbiorach nie jest globalnie jednostajna]
\label{ex:topog-odwzorowania-gorki-uciekajace}
Niech $X=Y=\R$ ze zwykłymi
topologiami i metryką na $Y$.
Dla $n\ge1$ określamy
\[
f_n(x)=\max\{1-|x-n|,0\}.
\]
Każda $f_n$ jest ciągła oraz
$\sup_{x\in\R}|f_n(x)|=1$.
Jeżeli $K\subseteq\R$ jest zwarty,
to jest ograniczony; dla
dostatecznie dużego $n$
mamy $f_n|_K=0$.
Zatem $f_n\to0$ jednostajnie
na każdym zbiorze zwartym,
lecz nie jednostajnie na całej $\R$.
\end{przyklad}

\section{Ewaluacja i prawo wykładnicze}
\label{sec:topog-odwzorowania-ewaluacja}

\begin{twierdzenie}[Ciągłość ewaluacji]
\label{thm:topog-odwzorowania-ewaluacja}
Niech $X$ będzie lokalnie zwartą
przestrzenią Hausdorffa, a $Y$
dowolną przestrzenią topologiczną.
Odwzorowanie ewaluacji
\[
\operatorname{ev}\colon
C_k(X,Y)\times X\to Y,\qquad
\operatorname{ev}(f,x)=f(x),
\]
jest ciągłe.
\end{twierdzenie}
\begin{proof}
Niech $(f_0,x_0)\in C_k(X,Y)\times X$
i niech $V\subseteq Y$ będzie otwartym
otoczeniem $f_0(x_0)$.
Zbiór $f_0^{-1}(V)$ jest otwartym
otoczeniem $x_0$ w $X$.
Z lematu o zwartych domknięciach
z rozdziału~14 istnieje
otwarty zbiór $W$ i zwarty zbiór $K$
takie, że
\[
x_0\in W\subseteq K
\subseteq f_0^{-1}(V).
\]
Wówczas $f_0\in[K,V]$,
a zbiór $[K,V]\times W$
jest otwartym otoczeniem $(f_0,x_0)$.
Dla dowolnego $(g,x)$ z tego
otoczenia mamy $x\in W\subseteq K$
i $g(K)\subseteq V$, więc
$\operatorname{ev}(g,x)=g(x)\in V$.
To wykazuje ciągłość w $(f_0,x_0)$.
\end{proof}

\begin{twierdzenie}[Prawo wykładnicze: funkcje z parametrem]
\label{thm:topog-odwzorowania-wykladnicze}
Niech $X$ będzie lokalnie zwartą
przestrzenią Hausdorffa, a $Z,Y$
dowolnymi przestrzeniami topologicznymi.
Niech $F\colon Z\times X\to Y$
będzie dowolnym odwzorowaniem.
Dla $z\in Z$ oznaczmy jego przekrój
przez $F_z\colon X\to Y$,
$F_z(x)=F(z,x)$.
Wówczas $F$ jest ciągłe
wtedy i tylko wtedy, gdy
każdy przekrój $F_z$ jest ciągły
oraz przyporządkowanie
$\widehat F\colon Z\to C_k(X,Y)$,
$\widehat F(z)=F_z$,
jest ciągłe.
W ten sposób przyporządkowanie
$F\mapsto\widehat F$ daje
bijekcję pomiędzy zbiorami
odwzorowań ciągłych
$C(Z\times X,Y)$ i
$C(Z,C_k(X,Y))$.
\end{twierdzenie}
\begin{proof}
Załóżmy, że $F$ jest ciągłe.
Dla każdego $z\in Z$
odwzorowanie $x\mapsto F(z,x)$
jest ciągłe jako złożenie
$F$ z ciągłym odwzorowaniem
$x\mapsto(z,x)$.
Niech teraz $[K,U]$ będzie
podbazowym zbiorem otwartym
w $C_k(X,Y)$ i niech
$z_0\in\widehat F^{-1}([K,U])$.
Wtedy $F(\{z_0\}\times K)\subseteq U$,
czyli otwarty zbiór $F^{-1}(U)$
zawiera $\{z_0\}\times K$.
Dla każdego $x\in K$ wybieramy
otwarte $O_x\ni z_0$ i $V_x\ni x$
z $O_x\times V_x\subseteq F^{-1}(U)$.
Zwartość $K$ daje skończone
$x_1,\ldots,x_m\in K$ takie, że
$K\subseteq\bigcup_{j=1}^mV_{x_j}$.
Wtedy $O=\bigcap_{j=1}^mO_{x_j}$
jest otwartym otoczeniem $z_0$
spełniającym $O\times K\subseteq
F^{-1}(U)$.
Jeśli $K=\varnothing$, bierzemy $O=Z$.
Zatem $O\subseteq
\widehat F^{-1}([K,U])$.
Wszystkie przeciwobrazy
zbiorów podbazowych są otwarte,
więc $\widehat F$ jest ciągłe.

Odwrotnie, załóżmy, że
$\widehat F\colon Z\to C_k(X,Y)$
jest ciągłe. Wówczas
\[
F=\operatorname{ev}\circ
(\widehat F\times\mathrm{id}_X).
\]
Pierwsze odwzorowanie w złożeniu
jest ciągłe jako iloczyn
odwzorowań ciągłych, a ewaluacja
jest ciągła na mocy poprzedniego
twierdzenia. Zatem $F$ jest ciągłe.
Obie operacje przyporządkowujące
sobie $F$ i $\widehat F$ są
wzajemnie odwrotne, co daje bijekcję.
\end{proof}

\begin{przyklad}[Dlaczego zwartość lokalna w ewaluacji nie jest zbędna]
\label{ex:topog-odwzorowania-ewaluacja-q}
Niech $X=\Q$ z topologią odziedziczoną
z $\R$, a $Y=\R$. Przestrzeń $\Q$
nie jest lokalnie zwarta i żaden
jej zwarty podzbiór nie ma
niepustego wnętrza w $\Q$
(patrz rozdział~14).
Wykażemy, że
$\operatorname{ev}\colon C_k(\Q,\R)
\times\Q\to\R$
nie jest ciągła w punkcie
$(0,0)$, gdzie $0$ oznacza
funkcję stale równą zero.

Niech $\mathcal N$ będzie dowolnym
otoczeniem zerowej funkcji w $C_k(\Q,\R)$,
a $V$ dowolnym otwartym otoczeniem
$0$ w $\Q$. Wybieramy bazowe
otoczenie
$\bigcap_{j=1}^m[K_j,U_j]
\subseteq\mathcal N$ funkcji zerowej,
gdzie $K_j\subseteq\Q$ są zwarte,
a $0\in U_j\subseteq\R$.
Zbiór $K=\bigcup_{j=1}^mK_j$
jest zwarty i ma puste wnętrze
w $\Q$; dla $m=0$ przyjmujemy
$K=\varnothing$. Istnieje więc
$r\in V\setminus K$.
Jeżeli $K\ne\varnothing$, funkcja
\[
g(q)=\min\left\{1,
 \frac{d(q,K)}{d(r,K)}\right\},
 \qquad q\in\Q,
\]
jest ciągła, zerowa na $K$
i spełnia $g(r)=1$.
Ponieważ $K$ jest domknięty
w przestrzeni metrycznej $\Q$
i $r\notin K$, mianownik jest dodatni.
Dla $K=\varnothing$ wybieramy
funkcję $g\equiv1$.
W obu przypadkach $g\in\mathcal N$,
ale $g(r)=1$. Stąd żadne
otoczenie punktu $(0,0)$ nie jest
odwzorowywane przez ewaluację
do $(-1/2,1/2)$, co przeczy
ciągłości w tym punkcie.
\end{przyklad}

\section{Metryzacja i zupełność przestrzeni odwzorowań}
\label{sec:topog-odwzorowania-metryzacja}

\begin{twierdzenie}[Metryka z wyczerpania zwartego]
\label{thm:topog-odwzorowania-metryka-wyczerpanie}
Niech $X$ będzie lokalnie zwartą,
$\sigma$-zwartą przestrzenią Hausdorffa,
a $(Y,d)$ przestrzenią metryczną.
Wybierzmy wyczerpanie zwarte
$(K_n)_{n\ge1}$ przestrzeni $X$
z rozdziału~14, tak że
$K_n\subseteq\operatorname{int}_XK_{n+1}$
oraz $X=\bigcup_{n\ge1}\operatorname{int}_XK_n$.
Dla $f,g\in C(X,Y)$ połóżmy
\[
\rho_n(f,g)=
\sup_{x\in K_n}\min\{1,d(f(x),g(x))\},
\qquad
D(f,g)=\sum_{n=1}^\infty
2^{-n}\rho_n(f,g).
\]
Przy $K_n=\varnothing$ ustalamy
$\rho_n(f,g)=0$.
Wówczas $D$ jest metryką zgodną
z topologią zwarto-otwartą na $C(X,Y)$.
\end{twierdzenie}
\begin{proof}
Każda $\rho_n$ jest nieujemna,
symetryczna i nie większa od $1$.
Nierówność trójkąta wynika
z nierówności
$\min\{1,a+b\}\le
\min\{1,a\}+\min\{1,b\}$
dla $a,b\ge0$, stosowanej punktowo,
a następnie z przejścia do supremum.
Dlatego szereg definiujący $D$
jest zbieżny i spełnia nierówność
trójkąta. Jeżeli $D(f,g)=0$,
to $\rho_n(f,g)=0$ dla każdego $n$.
Ponieważ zbiory $K_n$ pokrywają $X$,
otrzymujemy $f=g$.
Zatem $D$ jest metryką.

Dla każdego zwartego $K\subseteq X$
wyczerpanie z rozdziału~14 daje
$N$ z $K\subseteq K_N$.
W konsekwencji zbieżność
$\rho_N(f_\lambda,f)\to0$
wymusza jednostajną zbieżność
$f_\lambda$ do $f$ na $K$:
dla dowolnego $0<\varepsilon<1$
nierówność $\rho_N<\varepsilon$
oznacza $d(f_\lambda(x),f(x))
<\varepsilon$ dla wszystkich $x\in K$.
Odwrotnie, zbieżność jednostajna
na każdym $K_n$ implikuje
$\rho_n(f_\lambda,f)\to0$
dla każdego $n$.
Dla $\varepsilon>0$ wybieramy
$N$ tak, by
$\sum_{n>N}2^{-n}<\varepsilon/2$.
Następnie zbieżność skończenie
wielu $\rho_n$, $n\le N$,
pozwala zapewnić
$\sum_{n=1}^{N}2^{-n}\rho_n
<\varepsilon/2$.
Zatem $D(f_\lambda,f)<\varepsilon$.
Oba rozumowania obowiązują
dla sieci, nie tylko dla ciągów.
Na mocy twierdzenia
\ref{thm:topog-odwzorowania-zwarto-otwarta-jednostajna}
metryka $D$ wyznacza
topologię zwarto-otwartą.
\end{proof}

\begin{twierdzenie}[Zupełność przestrzeni odwzorowań]
\label{thm:topog-odwzorowania-zupelnosc}
Niech $X$ będzie lokalnie zwartą,
$\sigma$-zwartą przestrzenią Hausdorffa,
a $(Y,d)$ zupełną przestrzenią metryczną.
Dla dowolnego wyczerpania zwartego
$(K_n)$ metryka $D$ z poprzedniego
twierdzenia jest zupełna na $C(X,Y)$.
\end{twierdzenie}
\begin{proof}
Niech $(f_j)_{j\ge1}$ będzie ciągiem
Cauchy'ego w metryce $D$.
Dla każdego $n$ i $0<\varepsilon<1$
istnieje $j_0$ takie, że
\[
D(f_j,f_k)<2^{-n}\varepsilon
\qquad(j,k\ge j_0).
\]
Ponieważ $2^{-n}\rho_n(f_j,f_k)
\le D(f_j,f_k)$, otrzymujemy
$\rho_n(f_j,f_k)<\varepsilon$,
a więc
$d(f_j(x),f_k(x))<\varepsilon$
dla wszystkich $x\in K_n$.
Zatem $(f_j|_{K_n})$ jest
jednostajnie ciągiem Cauchy'ego.
Dla każdego $x\in X$ ciąg
$(f_j(x))$ jest więc ciągiem
Cauchy'ego w $Y$ i ma granicę
$f(x)\in Y$.

Wykażemy ilościowo jednostajną
zbieżność na $K_n$. Niech
$\eta>0$ będzie dowolna.
W rozumowaniu powyżej wybieramy
$0<\varepsilon<\min\{1,\eta\}$
i odpowiadający mu indeks $j_0$.
Dla $j\ge j_0$, $k\ge j_0$
oraz $x\in K_n$ mamy
$d(f_j(x),f_k(x))<\varepsilon$.
Przechodząc przy ustalonych $j,x$
do granicy $k\to\infty$,
otrzymujemy
$d(f_j(x),f(x))\le\varepsilon<\eta$.
Wybór $j_0$ nie zależy od $x\in K_n$,
więc $f_j\to f$ jednostajnie na $K_n$.
Każdy $x\in X$ należy do
$\operatorname{int}_X K_n$
dla pewnego $n$, więc $f$
jest ciągła w otoczeniu $x$
jako jednostajna granica funkcji
ciągłych (o wartościach w przestrzeni
metrycznej). W konsekwencji
$f\in C(X,Y)$.
Zbieżność jednostajna na wszystkich
$K_n$ implikuje $D(f_j,f)\to0$
na mocy poprzedniego twierdzenia.
Dowodzi to zupełności.
\end{proof}

\begin{przyklad}[Nieprzeliczalna przestrzeń dyskretna]
\label{ex:topog-odwzorowania-nieprzeliczalna}
Niech $X$ będzie nieprzeliczalną
przestrzenią dyskretną, a $Y=\R$.
Zwarte podzbiory $X$ są dokładnie
zbiorami skończonymi. Wobec tego
$C_k(X,\R)=\R^X$ z topologią
produktową, czyli topologią
zbieżności punktowej.

Pokażemy, że przestrzeń ta
nie spełnia pierwszego aksjomatu
przeliczalności w funkcji zerowej.
Gdyby $(U_n)_{n\ge1}$ było
przeliczalną bazą otoczeń,
dla każdego $n$ wybralibyśmy
podstawowe otoczenie produktowe
$V_n\subseteq U_n$ funkcji zerowej.
Zbiór współrzędnych ograniczanych
przez $V_n$ jest skończony;
oznaczmy go przez $F_n$.
Ponieważ $X$ jest nieprzeliczalna,
istnieje $x_0\in X\setminus
\bigcup_{n\ge1}F_n$.
Otoczenie
$W=\{f\in\R^X:|f(x_0)|<1\}$
funkcji zerowej nie zawiera
żadnego $V_n$: dla każdego $n$
można wybrać funkcję równą $0$
na $F_n$ i równą $2$ w $x_0$.
Zatem żadne $U_n$ nie zawiera się
w $W$, bo $V_n\subseteq U_n$.
Otrzymujemy sprzeczność.
W szczególności $C_k(X,\R)$
nie jest metryzowalna.
Pokazuje to, dlaczego założenia
o istnieniu przeliczalnego
wyczerpania zwartego są istotne.
\end{przyklad}

\begin{uwaga}[Związki z dalszą teorią]
Topologia zwarto-otwarta jest
naturalna przy badaniu przestrzeni
homotopii, odwzorowań i rodzin
parametryzowanych.
Zbieżność jednostajna na zbiorach
zwartych łączy topologię ogólną
z analizą funkcjonalną i teorią
przestrzeni funkcji.
W rozdziale~19 przejdziemy do
przestrzeni polskich i zbiorów
borelowskich, korzystając m.in.
z metryzacji i zupełności
udowodnionych tutaj.
\end{uwaga}
